Ten matched-pair simulations were designed to separate the hypotheses; their quantitative predictions (strength and fracture-mode class) were written to a hashed file before launch (\texttt{experiments/\allowbreak predictions/\allowbreak stage4\_\allowbreak predictions.json}) and are compared with the outcomes in Fig.~\ref{fig:discriminating_experiments}. Median absolute error of the strength predictions: 6.5\% (mean 10\%, maximum 30\%); fracture-mode class correct in 8 of 10 cases.

\paragraph{Pore shape at equal net section (H-A vs H-B).} Circular coarse pores whose diameter equals the side of the square pores (same minimum solid fraction 0.51, lower porosity) reach 16.1\,N/m, statistically indistinguishable from the square mesh (17.0 in Stage~2, 17.0 in the re-run of the same design), while H-A's necking knockdown had predicted 11.3\,N/m; at equal porosity the round pores give 14.1\,N/m (predicted 9.8) because the net section falls to 0.50. The shape of coarse pores therefore does not matter; the knockdown observed for fine meshes is an absolute ligament-size effect (ligaments of 10--13\,\AA{} carry a net-section stress of 17--19\,N/m, ligaments of 18--24\,\AA{} 22--31\,N/m). H-A is revised accordingly (H-A$'$: net section times a size- and alignment-dependent net-section strength); H-B explains the size dependence but not the orientation effects.

\paragraph{Ligament dispersion (H-C).} Dispersing the pore sizes by $\pm15\%$ lowered the strength of the square coarse mesh from 17.0 to 15.8\,N/m without changing its single-avalanche mode, and lowered the 32\,\AA{} round-pore mesh from 16.4 to 13.8\,N/m while \emph{shortening} its post-peak tail (failure strain 0.26 $\to$ 0.13, localisation 0.12 $\to$ 0.46). H-C is refuted: making parallel load paths non-equivalent is not sufficient for progressive failure. The pattern is that of an equal-load-sharing fibre bundle: when a strip fails, its load is transferred to the remaining ones, and with only three wide strips (or with a weakest link that is much weaker than the rest) the transfer exceeds their margin and the failure cascades; fine meshes with 8--10 rows of small ligaments, and jittered meshes, stay below that threshold and fail sequentially.

\paragraph{Flaw-size plateau (H-E).} A 30\,\AA{} and a 40\,\AA{} crack in 150\,\AA{} cells give 19.5 and 17.9\,N/m (predicted 17--19; Griffith scaling from the 10 and 20\,\AA{} cracks would give 13.9 and 12.0), and a 30\,\AA{} hole gives 18.6\,N/m (20\,\AA{} hole: 19.1). The plateau is real and not a periodic-image artefact: in this discrete lattice the tip bond configuration and its trapping strength, not the flaw length, set the instability once the flaw exceeds a few nanometres.

\paragraph{Vein direction (H-F).} With the same porosity and fine level, veins only along the load give 16.1\,N/m and $W=1.6$\,J/m$^2$, veins in both directions 12.8\,N/m, and veins only across the load 12.0\,N/m -- the value of the fine mesh alone. Load-perpendicular veins are dead weight (they are not crack arrestors, H-F$'$ rejected), and the strength of a nested mesh is that of its finest level plus the contribution of the load-parallel coarse ligaments, as H-F predicted (14.5 and 11.5). The load-parallel-only design is also the only nested mesh that combines a strength close to the coarse-only limit with progressive failure.
